\chapter{The Valuation-Adjustment Desk}\label{ch:rc:the-valuation-adjustment-desk}

A salesperson prices a twenty-year swap for an unrated industrial company that wants to
fix the rate on a long loan. The rates desk quotes the par rate in seconds. Then the
valuation-adjustment desk adds its charge: the company's credit, the bank's funding of the
uncollateralised position, the capital the trade will tie up for two decades. In this
chapter's example the charge comes to nearly forty basis points a year, more than the
bid--offer spread of the swap itself, and the client may walk to a bank with a cheaper view
of any of the three. Chapters~17 to~19 built the adjustments; this chapter is about the desk
that owns them: why it is central, how it prices a new trade against everything the bank
already has with the client, what it can hedge and what it cannot, and how its charges reach
the trading desks.

\section{Why one central desk}

\begin{definition}[XVA desk]\label{def:rc:the-valuation-adjustment-desk:desk}
An \emph{XVA desk}\index{XVA desk} is the trading desk that owns a bank's valuation
adjustments: it charges each new trade its adjustments, holds the resulting counterparty,
funding and capital risks on its book, hedges what can be hedged, and reports the
adjustments' changes as its P\&L.
\end{definition}

Adjustments are properties of \omterm{def:rc:counterparty-exposure:netting}{netting sets}, not of trades: a swap sold by the rates desk
and an FX forward sold by the FX desk to the same client net against each other, share the
client's collateral and its default. Only a desk that sees the whole \omterm{def:rc:counterparty-exposure:netting}{netting set} can price a
trade's effect on it. Centralising also concentrates the expertise and the hedging, which
needs default swaps, index protection and market hedges across asset classes
(\cref{fig:rc:the-valuation-adjustment-desk:org}).

\begin{omfigure}
\begin{tikzpicture}[font=\small,box/.style={draw,rounded corners,minimum width=2.3cm,minimum height=0.85cm,align=center}]
\node[box] (r) at (0,1.2) {rates desk};
\node[box] (f) at (0,0) {FX desk};
\node[box] (e) at (0,-1.2) {other desks};
\node[box,fill=omDef!10,minimum height=1.6cm] (x) at (5,0) {XVA desk\\(netting sets,\\CVA, FVA, KVA)};
\node[box] (m) at (10,1.2) {credit market\\(CDS, index)};
\node[box] (h) at (10,0) {rates and FX\\hedges};
\node[box,fill=omThm!10] (t) at (10,-1.2) {treasury and\\capital};
\foreach \n in {r,f,e} {\draw[->,thick] (\n.east) -- (x.west);}
\node[above,align=center,font=\footnotesize] at (2.5,1.05) {trades, and pays\\the charge};
\draw[->,thick] (x.east) -- (m.west);
\draw[->,thick] (x.east) -- (h.west);
\draw[<->,thick] (x.east) -- (t.west);
\node[above,font=\footnotesize] at (7.5,1.05) {hedges};
\end{tikzpicture}

\omcaption{The \omterm{def:rc:the-valuation-adjustment-desk:desk}{XVA desk} between the trading desks and the markets: desks pass it the
counterparty, funding and capital risks of their trades and pay a charge; it hedges credit
and market risk in the market and settles funding and capital with the treasury.
Schematic.}\label{fig:rc:the-valuation-adjustment-desk:org}
\end{omfigure}

\section{Incremental pricing and allocation}

\begin{definition}[Incremental XVA]\label{def:rc:the-valuation-adjustment-desk:incremental}
The \emph{incremental XVA}\index{incremental XVA} of a new trade is the change in the
\omterm{def:rc:counterparty-exposure:netting}{netting set}'s (or the portfolio's) adjustments when the trade is added:
$\mathrm{XVA}(\text{set}+\text{trade})-\mathrm{XVA}(\text{set})$. It is what the trade costs the bank,
and it can be negative when the trade offsets existing exposure.
\end{definition}

\begin{example}[Two new swaps in one netting set]\label{ex:rc:the-valuation-adjustment-desk:incremental}
Chapter~17's \omterm{def:rc:counterparty-exposure:netting}{netting set} with chapter~18's counterparty has a CVA of USD~891\,342. A new
five-year USD~50 million swap on which the bank pays fixed has a standalone CVA of 37\,061,
but adding it lowers the set's CVA to 869\,582: its incremental CVA is $-\text{USD}~21\,761$, a
rebate. The same swap with the bank receiving fixed has a standalone CVA of 30\,742 and an
incremental CVA of $+22\,784$.
\end{example}

\omcode{code/firm/xvaquote/firm_xvaquote.py}{39}{50}{Euler allocation of a netting set's CVA to its trades.}\label{lst:rc:the-valuation-adjustment-desk:euler}

Incremental charges depend on the order of arrival: the first trade with a client pays the
full standalone charge, a later offsetting one earns a rebate. To report each trade's share
of an existing adjustment, the desk needs an allocation that adds up.

\begin{definition}[Euler allocation]\label{def:rc:the-valuation-adjustment-desk:euler}
\emph{Euler allocation}\index{Euler allocation} assigns to trade $i$ the contribution
$a_i = \frac{d}{d\epsilon}\mathrm{XVA}(\ldots,(1+\epsilon)V_i,\ldots)|_{\epsilon=0}$. When the adjustment is
homogeneous of degree one in the trades' sizes, as CVA is, Euler's theorem makes the
contributions add up to the total; for CVA,
$a_i = (1-R)\sum_k\E[D(0,t_k)V_i(t_k)\mathbf 1_{V(t_k)>0}]\,\Delta\mathrm{PD}_k$.
\end{definition}

\begin{example}[Allocating the netting set]\label{ex:rc:the-valuation-adjustment-desk:euler}
Of the set's USD~891\,342 of CVA, \omterm{def:rc:the-valuation-adjustment-desk:euler}{Euler allocation} assigns 160\,167 to the swap and 731\,175 to
the cross-currency swap, against standalone CVAs of 288\,250 and 768\,334
(\cref{fig:rc:the-valuation-adjustment-desk:euler}). The swap, which offsets part of the
cross-currency exposure, receives most of the netting benefit.
\end{example}

\begin{omfigure}
\begin{tikzpicture}
\begin{axis}[width=11.5cm,height=5.2cm,ybar,bar width=24pt,
  symbolic x coords={swap,cross-currency},xtick=data,enlarge x limits=0.5,
  ylabel={CVA (USD thousand)},
  tick label style={font=\small},label style={font=\small},
  ymin=0,ymax=900,grid=major,grid style={gray!25},
  nodes near coords,every node near coord/.append style={font=\footnotesize,/pgf/number format/fixed,/pgf/number format/precision=0},
  legend columns=2,legend style={font=\footnotesize,at={(0.5,-0.2)},anchor=north,draw=none,fill=none,/tikz/every even column/.append style={column sep=8pt}},
  table/col sep=comma]
\addplot[fill=omExo!40,draw=omExo] table[x=trade,y=standalone]{figdata/rates-credit-risk/20-the-valuation-adjustment-desk/euler.csv};
\addplot[fill=omDef!60,draw=omDef] table[x=trade,y=euler]{figdata/rates-credit-risk/20-the-valuation-adjustment-desk/euler.csv};
\legend{standalone,Euler contribution}
\end{axis}
\end{tikzpicture}

\omcaption{Standalone CVA of each trade against its Euler contribution to the \omterm{def:rc:counterparty-exposure:netting}{netting set}'s
CVA. The contributions add up to the set's CVA; the difference from the standalone figures
is the netting benefit, shared in proportion to each trade's role in the exposure. Data:
the chapter's tutorial.}\label{fig:rc:the-valuation-adjustment-desk:euler}
\end{omfigure}

\section{What it hedges and what it cannot}

CVA has two kinds of risk: the market factors that move the exposure, hedged with swaps,
options and FX like any derivative book, and the counterparty's credit, hedged with credit
instruments. Many counterparties have no traded default swaps: the desk needs a spread to
price them and a proxy to hedge them.

\begin{definition}[Proxy credit spread]\label{def:rc:the-valuation-adjustment-desk:proxy}
A \emph{proxy credit spread}\index{proxy credit spread} is a spread assigned to a
counterparty without liquid credit quotes, estimated from liquid peers that resemble it in
credit quality, sector and region, or mapped to a single related name. In the European Union a proxy
used for CVA capital must consider all three attributes, with at least three industry categories (public,
financial, other) and four regions (Europe, North America, Asia, rest of the world).
\end{definition}

\begin{example}[A proxy for the unrated client]\label{ex:rc:the-valuation-adjustment-desk:proxy}
Regress the logarithm of eighteen illustrative ten-year peer spreads on rating (A, BBB, BB),
sector (industrials, utilities, consumer) and region (US, EU) indicators. The fit's error is
1.3\% of the spread, and for the client, mapped internally to BB, industrials, US, it predicts
337.5 basis points at ten years (\cref{fig:rc:the-valuation-adjustment-desk:proxy}). Applying
chapter~18's term-structure shape gives a curve from 154 basis points at one year to 338 at
ten.
\end{example}

\begin{omfigure}
\begin{tikzpicture}
\begin{axis}[width=8.5cm,height=6cm,
  xlabel={observed peer spread (bp)},ylabel={fitted spread (bp)},
  tick label style={font=\small},label style={font=\small},
  xmin=50,xmax=420,ymin=50,ymax=420,grid=major,grid style={gray!25},
  legend pos=north west,legend style={font=\footnotesize},
  table/col sep=comma]
\addplot[omExo,dashed,domain=50:420,samples=2,forget plot] {x};
\addplot[omMeth,only marks,mark=*,mark size=2pt] table[x=observed,y=fitted]{figdata/rates-credit-risk/20-the-valuation-adjustment-desk/proxy_a.csv};
\addplot[omDef,only marks,mark=square*,mark size=2pt] table[x=observed,y=fitted]{figdata/rates-credit-risk/20-the-valuation-adjustment-desk/proxy_bbb.csv};
\addplot[omThm,only marks,mark=triangle*,mark size=2.6pt] table[x=observed,y=fitted]{figdata/rates-credit-risk/20-the-valuation-adjustment-desk/proxy_bb.csv};
\legend{A peers,BBB peers,BB peers}
\end{axis}
\end{tikzpicture}

\omcaption{Observed against fitted ten-year spreads of the illustrative peers in the
proxy regression on rating, sector and region. The points sit close to the diagonal; the
client's proxy is the model's prediction for its cell (BB, industrials, US). Data: the
chapter's tutorial.}\label{fig:rc:the-valuation-adjustment-desk:proxy}
\end{omfigure}

\begin{definition}[Contingent credit default swap]\label{def:rc:the-valuation-adjustment-desk:ccds}
A \emph{contingent credit default swap}\index{contingent credit default swap} pays, on the
default of a reference entity, the \omterm{def:rc:reduced-form-credit:pd}{loss given default} on the then value of a specified
derivative (or \omterm{def:rc:counterparty-exposure:netting}{netting set}) instead of on a fixed notional: exact protection for the CVA of
that derivative, and priced by its CVA.
\end{definition}

A proxy hedge leaves basis risk (the client's credit does not move with its peers),
\omterm{def:rc:reduced-form-credit:jtd}{jump-to-default risk} (index protection does not pay when the client defaults alone) and
\omterm{def:rc:rates-risk:crossgamma}{cross-gamma} (exposure and credit move together). Funding and capital charges have no market
hedge at all: the desk manages them by collateral, compression, novation and the choice of
counterparties.

\begin{dated}{2026-09}{Proxy spreads and eligible hedges}\label{dat:rc:the-valuation-adjustment-desk:proxy}
The Basel Committee's CVA framework (July 2020) requires proxy spreads for illiquid
counterparties to be estimated from liquid peers by an algorithm that discriminates at least
on credit quality, industry and region, allows mapping to one related name (a municipality
to its country) only with supervisory justification, and recognises as CVA hedges only
single-name default swaps, single-name contingent default swaps and index default swaps.
\end{dated}

\section{Charging, transfer pricing and the desk's P\&L}

\begin{definition}[XVA charge]\label{def:rc:the-valuation-adjustment-desk:charge}
The \emph{XVA charge}\index{XVA charge} is the amount the \omterm{def:rc:the-valuation-adjustment-desk:desk}{XVA desk} takes from the
originating desk when a trade is booked: the trade's incremental adjustments, paid upfront or
converted into a running spread on the trade, $\text{charge} = \mathrm{XVA}/(N\cdot A)$ with $A$ the
trade's annuity.
\end{definition}

\begin{example}[The twenty-year quote]\label{ex:rc:the-valuation-adjustment-desk:quote}
The client receives floating and pays the par rate of 4.11\% on USD~100 million for twenty
years; no CSA. On 4\,000 paths the bank's \omterm{def:rc:counterparty-exposure:profiles}{expected exposure} peaks at USD~5.65 million after
five years (\cref{fig:rc:the-valuation-adjustment-desk:ee}). With the proxy curve, the CVA is
USD~1\,649\,494; funding at 80 basis points costs 343\,872; with capital charged at a 10\% hurdle
(SA-CCR at a 100\% risk weight, and CVA capital at the 7\% risk weight for unrated
industrials), the KVA is 3\,141\,826. Over the swap's annuity of 13.67, that is 12.1, 2.5 and 23.0
basis points a year: 37.6 in all.
\end{example}

\begin{omfigure}
\begin{tikzpicture}
\begin{axis}[width=11.5cm,height=5.2cm,
  xlabel={time (years)},ylabel={EE (USD million)},
  tick label style={font=\small},label style={font=\small},
  xmin=0,xmax=20,ymin=0,ymax=6.5,grid=major,grid style={gray!25},
  table/col sep=comma]
\addplot[omDef,very thick] table[x=t,y=ee]{figdata/rates-credit-risk/20-the-valuation-adjustment-desk/ee20.csv};
\end{axis}
\end{tikzpicture}

\omcaption{\omterm{def:rc:counterparty-exposure:profiles}{Expected exposure} of the twenty-year swap on which the bank receives fixed from
the unrated client: it peaks after about five years and runs off over the remaining fifteen.
Data: the chapter's tutorial.}\label{fig:rc:the-valuation-adjustment-desk:ee}
\end{omfigure}

The desk's P\&L is the change in the adjustments it holds, net of its hedges and of the
charges it receives. Charges set at inception and adjustments marked daily diverge: spreads
move, exposures grow, the bank's own funding cost changes. A desk that charged too little
loses slowly; one that charges too much loses the trades.

\section{Organisation}

Where the desk sits decides what it optimises. Inside the markets division it prices and
hedges like a trading desk; inside treasury or finance it is closer to a utility setting
transfer prices. Either way it needs the exposure engine of chapter~17 on common scenarios
for every \omterm{def:rc:counterparty-exposure:netting}{netting set}, legal data on netting and collateral agreements, credit data on every
counterparty, and a capital calculator, all fast enough to answer a salesperson in seconds:
the risk engine of chapter~29.

\section{Tutorial: pricing a new trade}\label{tut:rc:the-valuation-adjustment-desk:tutorial}

\begin{tutorial}
\textbf{Goal.} Compute incremental and allocated CVA in an existing \omterm{def:rc:counterparty-exposure:netting}{netting set}, build a proxy
spread, and quote the adjustments of a new long-dated swap as running basis points.
\textbf{End state:} the numbers of
\cref{ex:rc:the-valuation-adjustment-desk:incremental,ex:rc:the-valuation-adjustment-desk:euler,ex:rc:the-valuation-adjustment-desk:proxy,ex:rc:the-valuation-adjustment-desk:quote}
and the charts.
\begin{tutsteps}
\item \textbf{Incremental}: \texttt{incremental\_table()}, standalone against incremental.
\item \textbf{Allocation}: \texttt{allocation()}; check that the contributions add up.
\item \textbf{Proxy}: \texttt{proxy()} and \texttt{proxy\_curve()}.
\item \textbf{Quote}: \texttt{twenty\_year\_quote()}; \texttt{fig\_rc\_xvadesk.py} writes the charts.
\end{tutsteps}
\textbf{What to change next.} Map the client to BBB instead and see the quote move; add a CSA
and recompute every component.
\end{tutorial}

\section{Build: the quote service}\label{bld:rc:the-valuation-adjustment-desk:xvaquote}

\begin{build}
\bfield{Purpose} The firm's \omterm{def:rc:credit-and-debit-valuation-adjustments:xva}{XVA} quotes: incremental adjustments of a new trade against the
existing \omterm{def:rc:counterparty-exposure:netting}{netting set}, allocations for reporting, proxy spreads, and running charges.
\bfield{Interface} \texttt{netting\_cva(values, D, times, cpty, recovery)};
\texttt{incremental(adjustment, existing, new)}; \texttt{euler\_cva}; \texttt{proxy\_spread(peers,
target)}; \texttt{running\_charge(upfront, notional, annuity)}. Reads \texttt{firm\_exposure},
\texttt{firm\_cva}, \texttt{firm\_xvafund}.
\bfield{Rules} All adjustments on the same scenarios as the existing \omterm{def:rc:counterparty-exposure:netting}{netting set}; proxies
regress log spreads on credit quality, sector and region.
\bfield{Acceptance tests} \texttt{code/firm/xvaquote/tests/}: Euler contributions add up and
scale with the trade; the incremental identity; the proxy regression recovers a
multiplicative table exactly; the running-charge conversion.
\bfield{Stretch} Incremental FVA and KVA on the portfolio; allocation of capital by Euler on
the SA-CCR formula; a proxy with liquidity and seniority factors; sub-second incremental runs
by storing path values.
\end{build}

\begin{omsources}
\item Basel Committee on Banking Supervision, \emph{Targeted revisions to the credit valuation
adjustment risk framework}, July 2020, MAR50.
\end{omsources}

\section{Exercises}

\begin{exercise}[$\star$]\label{exo:rc:the-valuation-adjustment-desk:1}
Why can a trade's incremental CVA be negative while its standalone CVA is positive?
\end{exercise}

\begin{exercise}[$\star$]\label{exo:rc:the-valuation-adjustment-desk:2}
Convert an upfront charge of USD~1 million into running basis points on a USD~100 million
twenty-year swap with the chapter's annuity.
\end{exercise}

\begin{exercise}[$\star$]\label{exo:rc:the-valuation-adjustment-desk:3}
What are the three variables a proxy spread must discriminate on under the Basel CVA
framework?
\end{exercise}

\begin{exercise}[$\star\star$]\label{exo:rc:the-valuation-adjustment-desk:4}
Show that the Euler contributions to CVA add up to the \omterm{def:rc:counterparty-exposure:netting}{netting set}'s CVA.
\end{exercise}

\begin{exercise}[$\star\star$]\label{exo:rc:the-valuation-adjustment-desk:5}
Why does the order of arrival of trades matter for incremental charges, and why not for
\omterm{def:rc:the-valuation-adjustment-desk:euler}{Euler allocations}?
\end{exercise}

\begin{exercise}[$\star\star$]\label{exo:rc:the-valuation-adjustment-desk:6}
Which risks does index protection leave open when it hedges the CVA of an unrated client?
\end{exercise}

\begin{exercise}[$\star\star\star$]\label{exo:rc:the-valuation-adjustment-desk:7}
\emph{Coding.} Compute the \omterm{def:rc:reduced-form-credit:cs01}{CS01} of the twenty-year swap's CVA on the proxy curve and the
notional of ten-year protection on a peer that hedges it.
\end{exercise}

\begin{exercise}[$\star\star\star$]\label{exo:rc:the-valuation-adjustment-desk:8}
\emph{Find the flaw.} ``Our \omterm{def:rc:the-valuation-adjustment-desk:charge}{XVA charge} covers the expected loss on the client, so the \omterm{def:rc:the-valuation-adjustment-desk:desk}{XVA
desk}'s P\&L should be flat over the life of the trade.''
\end{exercise}

\section{Problem: The Unrated Corporate}

\begin{problem}[{Weekend problem --- the all-in add-on}]\label{pb:rc:the-valuation-adjustment-desk:1}
The company of \cref{ex:rc:the-valuation-adjustment-desk:quote} asks for the twenty-year swap.
Another bank has quoted it 20 basis points below this bank's all-in rate.

\medskip\noindent\textbf{Part I --- The quote.}
\begin{enumerate}
\item Give the swap's par rate and its annuity.
\item Give the proxy spread at ten years and the internal mapping behind it.
\item Give CVA, FCA and KVA upfront and in running basis points.
\item Give the all-in add-on.
\item Which component dominates, and why for this client?
\end{enumerate}

\medskip\noindent\textbf{Part II --- Hedging.}
\begin{enumerate}[resume]
\item Which part of the add-on can the desk hedge in the market, as a share of the total?
\item Give the CVA's \omterm{def:rc:reduced-form-credit:cs01}{CS01} and the hedge in ten-year peer protection.
\item What does the hedge leave open?
\item Would a contingent default swap be available, and what would it cost?
\item How would the desk hedge the exposure's rate risk?
\end{enumerate}

\medskip\noindent\textbf{Part III --- The competition.}
\begin{enumerate}[resume]
\item Which assumptions could make the other bank 20 basis points cheaper?
\item Map the client to BBB instead: what would change?
\item Which terms could the bank offer to close the gap?
\item Why might the other bank be wrong?
\item Is walking away a loss for the bank?
\end{enumerate}

\medskip\noindent\textbf{Part IV --- Judgement.}
\begin{enumerate}[resume]
\item Who should bear the KVA: the rates desk, the \omterm{def:rc:the-valuation-adjustment-desk:desk}{XVA desk} or the bank?
\item How should the charge be split with the salesperson's revenue?
\item What would you tell the client about the charge?
\item State the \emph{named result}: the all-in add-on in running basis points and the part the
desk can hedge.
\item In one sentence: what does an \omterm{def:rc:the-valuation-adjustment-desk:desk}{XVA desk} sell to the rest of the bank?
\end{enumerate}
\end{problem}

\section{Interview questions}

\begin{interviewq}[$\star$ \iqroles{trader, bank}]\label{iq:rc:the-valuation-adjustment-desk:1}
Why do banks centralise valuation adjustments in one desk?
\end{interviewq}

\begin{interviewq}[$\star\star$ \iqroles{researcher}]\label{iq:rc:the-valuation-adjustment-desk:2}
What is \omterm{def:rc:the-valuation-adjustment-desk:euler}{Euler allocation}, and when does it add up?
\end{interviewq}

\begin{interviewq}[$\star\star$ \iqroles{trader}]\label{iq:rc:the-valuation-adjustment-desk:3}
How do you price the CVA of a client with no traded credit?
\end{interviewq}

\begin{interviewq}[$\star\star$ \iqroles{developer}]\label{iq:rc:the-valuation-adjustment-desk:4}
A salesperson needs an \omterm{def:rc:the-valuation-adjustment-desk:incremental}{incremental XVA} quote in two seconds. How do you design the service?
\end{interviewq}

\begin{interviewq}[$\star\star\star$ \iqroles{trader, risk}]\label{iq:rc:the-valuation-adjustment-desk:5}
What drives the \omterm{def:rc:the-valuation-adjustment-desk:desk}{XVA desk}'s P\&L, and what can go wrong in a crisis?
\end{interviewq}

\begin{interviewq}[$\star\star\star$ \iqroles{bank}]\label{iq:rc:the-valuation-adjustment-desk:6}
Should \omterm{def:rc:the-valuation-adjustment-desk:charge}{XVA charges} be passed to clients in full? Discuss competition and adverse selection.
\end{interviewq}
